\section{Introduction}
\label{sec:intro}

Reinforcement learning from verifiable rewards (RLVR) has emerged as a powerful paradigm for improving mathematical reasoning in large language models \citep{shao2024deepseekmath, guo2025deepseek}.
Unlike RLHF, which requires a learned reward model, RLVR uses deterministic answer verification---a response is correct (reward 1) or incorrect (reward 0).
Group Relative Policy Optimization \citep[GRPO;][]{shao2024deepseekmath} is a popular RLVR algorithm that avoids per-token value estimation by normalizing advantages within a group of $G$ sampled responses to the same prompt.

\paragraph{The small-group binary degeneracy.}
While GRPO works well at large group sizes ($G \geq 8$) with rich reward distributions, practical constraints often force small $G$.
Under \emph{sparse binary rewards} at $G{=}4$, a critical failure mode emerges: if all four responses are correct or all four are incorrect---which is common for easy or hard prompts---the group-mean-centered advantage $(r_i - \bar{r}) / \sigma$ collapses to exactly zero, producing no gradient.
We show empirically that Dr.GRPO \citep{liu2025understanding}, which fixes GRPO's length bias but retains mean-centered advantage, achieves only 27.8\% on the full GSM8K test set at $G{=}4$ with binary reward---barely above the 25.5\% base model (\Cref{tab:main}).

\paragraph{Our approach.}
We propose two complementary mechanisms to restore effective learning under sparse binary rewards at small $G$:

\begin{enumerate}[leftmargin=*,topsep=2pt,itemsep=1pt]
\item \textbf{Threshold-Anchored Signed Advantage (TASA).}
Instead of centering advantages around the group mean, we anchor them to a fixed threshold $c$ (set to 0.5 for binary rewards) and compute:
\[
A_i = \frac{(r_i - c)_+}{\Zp} - \frac{(c - r_i)_+}{\Zm},
\]
where $\Zp$ and $\Zm$ normalize the positive and negative excess, respectively.
This guarantees that every correct response receives a strictly positive advantage and every incorrect response receives a strictly negative one, regardless of the group composition.
TASA alone improves GSM8K accuracy to 46.8\%, a +19.0~pp gain over Dr.GRPO (\Cref{tab:main}).

\item \textbf{Verified CE Replay.}
We maintain a per-prompt, hash-deduplicated bank of verified successful completions (at most $K{=}2$ per prompt) and replay them during training with a standard supervised cross-entropy (CE) loss.
Unlike RePO \citep{li2025repo} and RLEP \citep{zhang2025rlep}, which replay with importance-weighted policy-gradient objectives requiring stored log-probabilities, our CE replay needs only token IDs and avoids importance-sampling staleness entirely.
Adding CE replay to TASA yields 55.0\%, a further +8.3~pp improvement (\Cref{tab:main}).
\end{enumerate}

\paragraph{Contrastive replay is suboptimal.}
As an ablation, we also test a DPO-style \citep{rafailov2023direct} contrastive pair replay that uses both verified successes and failures in a pairwise log-sigmoid objective.
This variant achieves only 49.0\%, underperforming positive CE replay by 6.0~pp.
The result suggests that, in sparse binary RLVR, directly imitating verified successes is more effective than contrasting them with failures---the negative signal may suppress useful reasoning steps that happen to appear in incorrect responses.

\paragraph{Contributions.}
\begin{itemize}[leftmargin=*,topsep=2pt,itemsep=1pt]
\item We identify and characterize the \emph{small-group binary degeneracy} in GRPO, where mean-centered advantage produces zero gradients for all-correct or all-incorrect groups (\Cref{sec:prelim}).
\item We propose TASA, a threshold-anchored signed advantage that provides non-degenerate gradients under binary reward. While the formulation is simple---a count-normalized sign update around a fixed threshold---it yields a +19~pp gain over Dr.GRPO at $G{=}4$, demonstrating that even straightforward fixes to the advantage can have large impact (\Cref{sec:tasa}).
\item We combine TASA with verified CE replay---a simpler alternative to the importance-weighted PG replay used by RePO \citep{li2025repo} and RLEP \citep{zhang2025rlep}---that further improves accuracy by +8.3~pp without requiring stored log-probabilities (\Cref{sec:replay}).
\item We show through ablation that, in our setting, positive CE replay outperforms a DPO-style contrastive pair replay by 6.0~pp, suggesting that direct success imitation is more effective than pairwise negative contrast under sparse binary verification (\Cref{sec:ablation}).
\end{itemize}
